\documentclass[twocolumn]{aastex631}
%\usepackage{draftwatermark}

\graphicspath{{/Users/devaldeliwala/research/muLab/paper/plots/}}

%\SetWatermarkText{DRAFT}
%\SetWatermarkScale{10}
%\SetWatermarkColor[gray]{0.97}
\usepackage{float} % for 'H' placement
\usepackage{multirow}
\usepackage{natbib}


\begin{document}

\title{THE NEAR--INFRARED DISCRETE EXTINCTION LAWS OF THE GALACTIC CENTER}

\author{Deval Deliwala}
\affiliation{University of California, Berkeley}

\author{Matthew Hosek} 
\affiliation{University of California, Los Angeles} 

\author{Jessica Lu} 
\affiliation{University of California, Berkeley} 

\shorttitle{Extinction Laws of Galactic Center}
\shortauthors{Deliwala, Hosek, Lu, $\cdots$ }

\begin{abstract}
We investigate the discretized extinction law in the Galactic Center (GC) by
applying the Red Clump (RC) method to JWST NIRCam catalogs. Analyzing four
filters -- 1.15$\mu$m, 2.12$\mu$m, 3.23$\mu$m, and 4.05$\mu$m --
we derived extinction ratios, $ A_\lambda/A_{K_s}$, where $K_s \approx
2.12\mu$m, across four spatial regions NRCB1--4 covering the
central 120"$\times$120" of the GC.  To enhance the accuracy
of the RC bar slope determination,
we developed a novel Optimized Curve Fitting (OCF) algorithm. Our results reveal
significant spatial variations in extinction ratios among the regions. At 1.15$\mu$m 
and 4.05$\mu$m, the extinction ratios across all regions converge
towards similar values consistent with previous studies. However, at 3.23$\mu$m, 
we find significant deviations between each region, with NRCB3 exhibiting the highest
extinction ratio of $4.18 \pm 0.43$, and NRCB2 the lowest at $2.15 \pm 0.37$. 
While the extinction ratios $3.23\mu$m for NRCB1, NRCB2, and NRCB4 are
not statistically different from each other, NRCB3 is
significantly higher and statistically different from NRCB2 and the
\citep{Fritz2011} law, which covers the inner $14"\times 20"$ of the NRCB3 region.
These findings underscore spatial variations in extinction
within the Galactic Center, particularly in the relative vicinity of Sag. A* in
NRCB3. The observed variations suggest that local physical conditions—such as ice
absorption features—likely contribute to the extinction differences.
\end{abstract}

\section{Introduction}

Understanding the interstellar extinction law is crucial for studying the
structure and star formation history of the Galactic Center (GC) and other dense
regions within the Milky Way. Interstellar extinction refers to the absorption
and scattering of light by dust grains along the line of sight, leading to the
reddening of stellar objects relative to their intrinsic magnitudes \citep[e.g.][]{Draine2003}. The extinction law, which describes how this effect
varies with wavelength, is essential for correcting astronomical observations
and accurately interpreting the physical characteristics of stars and other
objects in these obscured regions. Common extinction laws like \citep{Cardelli1989}
and \citep{Fitzpatrick1999} describe the optical infrared (OIR) extinction as a single power law
($A_\lambda \propto \lambda^{-\beta}$) with $\beta \approx 1.6$.  More recent studies
\citep[e.g.][]{Nishiyama2009,Fritz2011}
have used large-scale infrared surveys  (e.g. IRSF/SIRIUS, 2MASS)
to improve the measurement of the OIR extinction law.
A power law has continued to be found a good fit in the NIR regime, though with a steeper
exponent than what was found in previous works.

The extinction law is not uniform throughout the Galaxy; it depends on both the
line of sight and the physical properties of the intervening dust grains
\citep{Fitzpatrick2007}. In regions such as the GC, where
extreme dust obscuration occurs, determining an accurate extinction law is even
more critical to correct for these effects and to understand the underlying
stellar populations \citep{Nishiyama2009,NoguerasLara2018}.

One of the foundational techniques for deriving the extinction law in dense
regions is the ``Red Clump (RC) method," introduced by \citep{Wozniak1996}.
RC stars are a population of evolved stars in their core-helium burning phase,
serving as excellent standard candles due to their relatively uniform intrinsic
luminosities and colors \citep{Girardi2001}. By analyzing the mean
magnitudes and colors of red giant stars (often RC  stars) across different wavelength bands,
\citep{Messineo2005,Nishiyama2009} reported steeper
extinctions of $\beta = 1.9$, and 2.0 respectively towards the Galactic Bulge.

\citep{Fritz2011} further contributed by deriving the central GC extinction law
through spectroscopic observations of gas emission lines ($\lambda \geq 1\mu$m) within 20 arcseconds of
Sagittarius A*. By analyzing hydrogen emission lines and stellar colors, they
measured the wavelength-dependent extinction and described it using a power-law
with $\beta = 2.11 \pm 0.06$.

Advancements in infrared astronomy and the availability of deep photometric
catalogs have  facilitated the study of extinction laws. For instance,
\citep{NoguerasLara2018,NoguerasLara2020} conducted high-angular-resolution ($\sim$0.2'')
\textit{JHKs} imaging analysis of the GC using the GALACTICNUCLEUS survey which covers
a range of $\sim$6000pc$^2$.  Similarly, studies utilizing data from the VISTA \textit{Variables in the Via
Lactea} (VVV) survey \citep{Minniti2010} and the Galactic Legacy
Infrared Mid-Plane Survey Extraordinaire (GLIMPSE) survey
\citep[e.g.][]{Benjamin2003} have
provided extensive datasets for mapping interstellar extinction across the
Galactic plane. However, survey-based extinction studies are often dominated by low-extinction
stars and are limited by the photometric depths of the survey(s).

\citep{Hosek2018} presented an updated GC extinction
law in the optical and near-infrared regime (0.8$\mu$m -- $2.2\mu$m) using Hubble Space Telescope (HST) observations of 819 RC stars. This work found the GC extinction law to be inconsistent with a single power law and found extinction ratios at 1.25$\mu$m and
$8.14\mu$m relative to the $Ks$ (2.12$\mu$m) band to be 18\% and 24\% larger than the previously
used \citep{Nishiyama2009} law.

The recent release of photometrically calibrated catalogs from the James Webb
Space Telescope (JWST) offers an unprecedented opportunity to probe the
extinction law in greater detail, both in terms of wavelength and possible spatial variation.

In this study, we build upon these foundational works by applying the RC method
to the new JWST NIRCam catalogs, which include four different filters spanning
from the near infrared (NIR) (F115W at 1.15$\mu$m, F212N at 2.12$\mu$m) to the
short-wavelength infrared (SWIR) (F323N at 3.23$\mu$m, F405N at 4.05$\mu$m).
We analyze four spatial regions near the Galactic Center—NRCB1, NRCB2, NRCB3,
and NRCB4—corresponding to different detectors on the JWST's sensor chip
assembly (see Figure \ref{fig:nrcb}). This approach allows for a more refined
determination of the extinction law, with particular focus on potential spatial
variations and the influence of physical phenomena such as ice absorption
features at specific wavelengths \citep[e.g.][]{Fritz2011}.

Our goal is to investigate whether significant differences in extinction law
behavior exist across these regions and to compare our findings with extinction
laws derived in past studies. By leveraging the high spatial resolution and
sensitivity of JWST data, we aim to enhance our understanding of the dust
properties in the GC and their impact on interstellar extinction. This work has
implications not only for studies of the GC but also for understanding dust
properties and extinction laws in other dense and obscured regions of the Milky
Way.

\section{Observations and Measurements}

% Information on Photometric Calibration of Catalogs
% What each NRCB region covers, how many stars in each region, filter, ...

\section{Methods}
\begin{figure}[t]
    \centering
    \includegraphics[width = \linewidth]{nrcb.pdf}
    \caption{JWST NRCB region coverage. }
    \label{fig:nrcb}
\end{figure}

The four NRCB regions of interest are labeled in Figure \ref{fig:nrcb}.
 To measure the ratios of total to selective extinction, such as $A_{J} / E_{J -
K_s}$ and $A_{Ks} / E_{J - Ks}$, we use the bulge RC stars, which make up a
well-defined linear clump in the GC color-magnitude diagrams (CMDs). The $J$
versus $J - K_s$ CMD constructed within the NRCB1 region, with 35,403 stars, is
shown in Figure \ref{fig:JVJ}, where the F115W (1.15 $\mu$m) filter corresponds
to the $J$ band and F212N (2.12 $\mu$m) corresponds to the $K_s$ band.

\begin{figure}[htp]
  \centering
    \includegraphics[width = \linewidth]{hess.png}
    \caption{NRCB1 $F115W - F212N$ vs. $F115W$ unsharp-masked CMD. }
    \label{fig:hess}
\end{figure}

\begin{figure}[htp]
  \centering
    \includegraphics[width = \linewidth]{JV-J.png}
    \caption{NRCB1 $F115W - F212N$ vs. $F115W$. The dashed parallel
    lines enclose the RC clump.}
    \label{fig:JVJ}
\end{figure}

The red dashed lines in Figure \ref{fig:hess} and \ref{fig:JVJ} enclose the RC stars. This cutoff
was determined by-eye using the unsharp-masking technique described in DeMarchi
et al. (2016), which emphasizes the high-density and de-emphasizes the low-density
features of the CMD. The $F115W - F212N$ vs. $F115W$ is converted into
a Hess Diagram, which represents stellar density as a 2D histogram. Each star's
position is modeled as a Gaussian distribution, with its width determined by the
associated photometric error. The histogram is constructed with a bin size of
0.05 mags in both color ($x$) and magnitude ($y$) space. To generate the
unsharp mask, we convolve the Hess Diagram with 2D Gaussian Kernel with a width
of 0.2 mags. Finally, the convolved Hess diagram is subtracted from the
original, yielding the unsharp mask shown in Figure \ref{fig:hess} which allows us to
roughly cutoff the RC bar shown in Figure \ref{fig:JVJ}.

The well-defined ``RC bar" is compact, with a slope that
defines the ratio of total
to selective extinction $A_{F115W} / E_{F115W - F212N}$. This slope
quantifies how far we have to ``push" stars across the CMD, in the direction
defined by the RC bar, to recover the intrinsic, un-reddened magnitudes of
stars. Calculating the RC slope is a non-trivial task, as it depends on how one
chooses to define the RC bar boundaries. To enhance the accuracy and minimize
human-induced errors, we developed an Optimized Curve Fitting (OCF) algorithm,
which requires minimal input parameters and outputs a weighted slope calculation
that is less dependent on initial cutoffs.

\subsection{Optimized Curve Fitting (OCF)}

The Optimized Curve Fitting (OCF) algorithm is designed to address the
irregularities and complexities inherent in CMD data, such as uneven data
distributions and background star contamination. Its primary objective is to fit
several compound Gaussian and linear models along the RC region, minimizing error metrics
across multiple iterations to accurately determine the net RC slope.

OCF begins by defining initial parameters and boundary conditions. These include
an estimate of the RC bar's extent, represented by two tight parallel lines
(shown as red dashed lines in Figure \ref{fig:JVJ} and \ref{fig:JVJcutoff}) that enclose the bar within
the CMD, and the range in color space over which the bar is analyzed. These
tight parallel cutoffs are not strictly enforced; they are expanded by a factor
of three (represented by the aqua lines in Figure \ref{fig:JVJcutoff}) to ensure
sufficient coverage.

\begin{figure}[t]
  \centering
    \includegraphics[width = \linewidth]{cutoff.png}
    \caption{NRCB1 $F115W - F212N$ vs. $F115W$. OCF input parameters
        for $n=15$.
        The aqua lines represent the 3x RC expansion. The red tiles represent
    the initial segmentation cutoff.}
    \label{fig:JVJcutoff}
\end{figure}

\subsubsection{Tile-Based Segmentation}

The RC bar is initially divided into $n$ tiles. For each tile, OCF attempts to
fit a compound model combining a Gaussian distribution (representing the RC
peak) and a linear component (representing the background trend). The width of
each tile is initialized to $x_{\text{range}} / n$, where $x_{\text{range}}$ is
the specified input color range of the RC bar. The height of each tile is
confined between the extended parallel lines, which are extrapolated by a factor
of three to prevent boundary conditions from interfering with the fitting
process. The tiles used as an initial ansatz for OCF are illustrated in Figure
\ref{fig:JVJcutoff}.

\subsubsection{Curve Fitting Procedure}

The core of the OCF algorithm is an iterative fitting process that employs
compound \verb|Gaussian1D + Linear1D| models from \verb|astropy.modeling| over
each tile. The algorithm first calculates a preliminary estimate of the RC
tile's mean and standard deviation using loose constraints imposed on the model
parameters, such as the amplitude and width of the Gaussian component, requiring
them to approximately fit the data distribution within the tile.

The preliminary fitting is iterated over a range of histogram bin counts
(typically 8-20). The resulting means and standard deviations from each fit are
averaged to provide a rough estimate of the RC peak parameters for each tile.
These initial fits serve as a foundation for a more refined fitting process,
where stricter bounds are placed on the parameters, especially the Gaussian
component's width and mean position.

\subsubsection{Error Minimization}

After obtaining the initial fits, the algorithm performs a second, more
stringent fitting pass to further refine the Gaussian parameters. The standard
deviation of the stringent fit is constrained within a small tolerance,
determined by
\[
    \sigma_{\text{range}} = \bar{\sigma} \pm \xi,
\]
where $\bar{\sigma}$ is the averaged standard deviation from the initial fits,
and $\xi$ is a small value (e.g., 0.2) to ensure the Gaussian part of the fit
matches the RC peak. The same process is then repeated with these more stringent
requirements.

To validate the quality of the fit, a synthetic cluster of stars is generated
using inverse transform sampling based on the cumulative distribution function
(CDF) of the compound fit. An example of a synthetic distribution is shown in
Figure \ref{fig:synthetic_dist}. A Kolmogorov-Smirnov (KS) test is employed
between the synthetic and observed data, which calculates the maximum difference
between the cumulative distribution functions of the two datasets. The fit is
considered successful if the test statistic is below a defined threshold
(typically 0.1).

\begin{figure}[htpb]
    \centering
    \includegraphics[width = \linewidth]{contour.pdf}
    \caption{NRCB1 $F115W - F212N$ vs. $F115W$. Actual vs. Synthetic
    Contour and Density Plots.}
    \label{fig:synthetic_dist}
\end{figure}

\subsubsection{Iteration Across Tiles}

This fitting process is applied to each tile along the RC bar. If the fitting
procedure fails in a given tile (e.g., due to insufficient data or poor
convergence), the tile width is slightly extended, and the process is repeated.
A maximum of seven tile extensions is allowed, which was empirically found to
optimize the balance between computational efficiency and accuracy. The start
position of the adjacent tile is then adjusted to the end of the previous tile.
The iteration completes once a tile's end position exceeds the provided
$x_{\text{range}}$ input parameter.

For each tile, the compound fit with the lowest error on the mean (EOM) under
the Gaussian component is selected as the optimal solution for that tile. The
result of the OCF procedure for NRCB1 F115W–F212N vs. F115W is shown in Figure
\ref{fig:ocfresult}. 

\subsubsection{Global Slope Determination}

Once the optimal fit parameters are determined for each tile, the global slope
of the RC bar is computed using a weighted \verb|scipy.optimize.curve_fit()|
method. The weights are derived from the inverse square of the EOM for each
tile, ensuring that tiles with more reliable fits contribute more significantly
to the slope measurement. Since the slope measurement may depend on the initial
number of ansatz tiles ($n$ parameter), OCF is repeated for $n = [15, 20]$. The
\textit{final} slope is chosen as the one with minimal error. A calculated final
slope of $1.367 \pm 0.007$ for NRCB1 F115W–F212N vs. F115W is shown in Figure
\ref{fig:slope}.

Unlike traditional methods that often require manual adjustments or heuristics
to define fitting boundaries and manage outliers, OCF automates these processes
through adaptive tiling and iterative error minimization. By dynamically
adjusting tile widths and applying stringent error metrics, OCF ensures a more
reliable and precise fit, even in regions with sparse data or high background
contamination. This automation reduces the risk of bias and allows for more
consistent analysis across diverse CMDs and filter combinations.

\begin{figure*}[htp]
  \centering
    \includegraphics[width = \linewidth]{results.pdf}
    \caption{NRCB1 $F115W - F212N$ vs. $F115W$. Each compound fit is
        optimized over all fit parameters. The RC bar is shown on the $xy$
    plane.}
    \label{fig:ocfresult}
\end{figure*}

\begin{table*}
\centering
\begin{tabular}{c||c|c|c|c}
$A_\lambda / E_{\lambda' - \lambda}$ & NRCB1 & NRCB2 & NRCB3 & NRCB4 \\[5px]
\hline
$A_{F115W} / E_{F115W - F212N}$ & $1.376 \pm 0.059$ & $1.369 \pm 0.043$ & $1.340 \pm 0.107$ & $1.404 \pm 0.006$ \\[5px]
$A_{F212N} / E_{F115W - F212N}$ & $0.400 \pm 0.050$ & $0.397 \pm 0.037$ & $0.378 \pm 0.037$ & $0.437 \pm 0.009$ \\[5px]
$A_{F212N} / E_{F212N - F323N}$ & $-0.635 \pm 0.107$ & $-0.793 \pm 0.107$ & $-0.343 \pm 0.030$ & $-0.509 \pm 0.031$ \\[5px]
$A_{F323N} / E_{F212N - F323N}$ & $-1.522 \pm 0.428$ & $-1.707 \pm 0.428$ & $-1.432 \pm 0.428$ & $-1.552 \pm 0.129$ \\[5px]
$A_{F212N} / E_{F212N - F405N}$ & $1.640 \pm 0.031$ & $1.505 \pm 0.026$ & $1.651 \pm 0.084$ & $1.437 \pm 0.028$ \\[5px]
$A_{F405N} / E_{F212N - F405N}$ & $0.656 \pm 0.030$ & $0.479 \pm 0.026$ & $0.397 \pm 0.098$ & $0.485 \pm 0.040$ \\[5px]
\end{tabular}
\caption{Calculated $A_\lambda / E_{\lambda' - \lambda}$ values for regions NRCB1–NRCB4. Filter labels replace wavelengths for clarity.}
\label{table:selective_extinction}
\end{table*}

\begin{table*}
\centering
\begin{tabular}{c||c|c|c|c||c}
$A_\lambda / A_{Ks}$ & NRCB1 & NRCB2 & NRCB3 & NRCB4 & \citep{Fritz2011} \\[5px]
\hline
$A_{F115W} / A_{Ks}$ & $3.44 \pm 0.06$ & $3.45 \pm 0.05$ & $3.54 \pm 0.11$ & $3.21 \pm 0.06$ & $3.044$ \\[5px]
$A_{F323N} / A_{Ks}$ & $2.40 \pm 0.43$ & $2.15 \pm 0.37$ & $4.18 \pm 0.43$ & $3.05 \pm 0.31$ & $0.64$ \\[5px]
$A_{F405N} / A_{Ks}$ & $0.40 \pm 0.03$ & $0.32 \pm 0.03$ & $0.24 \pm 0.06$ & $0.34 \pm 0.03$ & $0.38$ \\[5px]
\end{tabular}
\caption{Calculated $A_\lambda / A_{Ks}$ values for regions NRCB1–NRCB4 compared to \citep{Fritz2011}. Filter labels replace wavelengths for clarity.}
\label{table:total_extinction}
\end{table*}

\begin{figure}[ht]
  \centering
    \includegraphics[width = \linewidth]{slope_.png}
    \caption{NRCB1 $F115W - F212N$ vs. $F115W$. An RC slope of $1.367
    \pm 0.007$ was calculated using OCF.}
    \label{fig:slope}
\end{figure}

Another key advantage of the OCF algorithm is its ability to detect potential
subpopulations by segmenting the RC bar into different regions and independently
applying the fitting process to each. If the slopes derived from these regions
are statistically significantly different, it could indicate the presence of
subpopulations within the star cluster. This feature enables a deeper
investigation into the structure of stellar populations without relying on
pre-existing assumptions. However, as discussed later, no
statistically significant slope variations were found in this study, suggesting
that the RC is composed of a relatively uniform population within each NRCB region.
\section{Results}

The OCF procedure outlined above is performed for 6 filter combinations
spanning every $A_\lambda / E_{\lambda - \lambda'}$ for each NRCB regions
necessary to algebraically derive all $A_\lambda / A_{Ks}$.  Calculated OCF slopes for each filter-region combination
are shown in Table \ref{table:selective_extinction}. Calculated ratios of total extinction $A_\lambda / A_{Ks}$ are shown in
Table \ref{table:total_extinction} and are compared to values predicted by the \citep{Fritz2011} power law.

At $1.15$$\mu$m, the extinction ratios across all NRCB regions converge towards
a common value, with NRCB3 having the steepest extinction ratio ($A_{1.15\mu \rm
m} / A_{Ks}$) of $3.54 \pm 0.11$ and NRCB4 having the shallowest ratio of $3.21
\pm 0.06$. Statistical analysis indicates that the differences between NRCB1,
NRCB2, and NRCB3 are within their respective uncertainties, suggesting no
significant variation among these regions at this wavelength. The slightly lower
value in NRCB4 is statistically significant from NRCB2, but does not present a statistically
significant deviation when considering the combined uncertainties from all NRCB
regions. Across all regions, however, the calculated extinction ratios remain steeper than the
value of $3.044$ reported by \citep{Fritz2011}.

At $4.05$$\mu$m, the extinction ratios also show convergence toward a common
value, with NRCB1 having the steepest extinction ratio of $0.40 \pm 0.03$ and
NRCB3 having the shallowest value of $0.24 \pm 0.06$. Comparing these results
with the value reported by \citep{Fritz2011}, $0.38$, we find strong agreement
for NRCB1, NRCB2, and NRCB4. The value for NRCB3 is lower but remains within the
margin of error, indicating that the difference is also not statistically
significant.

The most notable variations emerge at the $3.23\mu$m wavelength, where we
observe significant differences in extinction ratios among the regions. NRCB3
exhibits the highest extinction ratio of $4.18 \pm 0.43$, significantly higher
than the other regions. NRCB2 has the lowest extinction ratio of $2.15 \pm
0.37$, while NRCB1 and NRCB4 have extinction ratios of $2.40 \pm 0.43$ and $3.05
\pm 0.31$, respectively. Statistical analysis shows that the extinction ratio in
NRCB3 is significantly different from NRCB2, whereas the differences
among NRCB1, NRCB2, and NRCB4 are within their uncertainties and not
statistically significant.

\begin{figure}[t]
  \centering
    \includegraphics[width = \linewidth]{slope_iso_.png}
    \caption{NRCB1 $F115W - F212N$ vs. $F115W$. The calculated
        extinction vector (red) compared to the \citep{Fritz2011} extinction
        vector (blue). The Fritz vector was generated by plotting the path
        traversed by stellar isochrones of increasing extinction using Fritz's
    power law. The isochrones were generated using SPISEA.}
    \label{fig:nr1result}
\end{figure}

\begin{figure}[htp]
  \centering
    \includegraphics[width = \linewidth]{extinction_law.png}
    \caption{Calculated discretized extinction law for each NRCB region. The
        black dashed line indicates the \citep{Fritz2011} law, calculated by
    analyzing the central $R<20''$ of the GC.}
\end{figure}

These variations at $3.23\,\mu$m suggest the presence of strong absorption
features, likely due to unevenly distributed ice absorption. The higher
extinction ratio in NRCB3, where Sagittarius A* is located, points to complex
dust and ice compositions in the vicinity of the supermassive black hole. The
observed deviations challenge the conventional understanding of interstellar
extinction laws in the Galactic Center.

The high uncertainties associated with the extinction ratios at $3.23\,\mu$m
stem from the complexities introduced by ice absorption features. The presence
of ice absorption complicates the extinction curve, making it difficult to model
using standard extinction laws. Accurately calculating the extinction slopes in
the presence of such features requires high-quality data and sophisticated
modeling techniques that can account for the additional absorption.

Given that our calculated extinction ratios at the $3.23\,\mu$m wavelength
differ significantly from previously established extinction laws like that of
\citep{Fritz2011}, it is essential to validate our methodology. A critical
parameter that greatly influences the final extinction ratios is the initial RC
cutoffs applied to the F323N CMDs within the OCF algorithm. Therefore, we must
ensure that the stars within these RC cutoffs are indeed RC stars to which the
OCF algorithm is fitting.

Several methods exist to verify the RC star selection, such as generating
synthetic star clusters or performing spectroscopic analyses. However, the most
straightforward approach is to cross-match stars from a well-defined RC cluster
in a different CMD onto the F323N CMDs. We utilize the F115W–F212N vs. F115W
CMD, shown previously for NRCB1, to define the RC stars. In this CMD, the RC bar
is prominently distinguished due to the substantial extinction difference
between the 1.15$\mu$m and 2.12$\mu$m wavelengths. This pronounced separation
allows for a clear identification of RC stars. We then apply a similar cutoff
strategy as initially employed in the OCF algorithm to isolate the RC population
accurately.

Subsequently the identified RC stars from the F115W–F212N versus F115W CMDs were
matched to their counterparts in the respective NRCB region F323N catalogs. Each star
was plotted with identical colors, enabling a direct comparison of how the RC distribution
changes across filters and allowing us to trace the RC slope consistently.

By following the same trajectory of RC stars' color transitions into the F323N filter, we found
that the extinction law from \citep{Fritz2011} does not align with the RC color distribution,
whereas the OCF-calculated extinction vector accurately follows it (see Figure
\ref{fig:nrcb3_plots}).

\section{Conclusion}

This study provides a detailed examination of the near-infrared extinction law
in the Galactic Center using JWST NIRCam data across four spatial regions:
NRCB1, NRCB2, NRCB3, and NRCB4. By employing the Red Clump (RC) method and a
novel Optimized Curve Fitting (OCF) algorithm, we calculated extinction ratios,
$A_\lambda / A_{Ks}$, for four filters: F115W, F212N, F323N, and F405N for every 
NRCB region. 

Our results demonstrate significant spatial variations in extinction ratios specifically 
at the 3.23$\mu$m wavelength. At
wavelengths 1.15$\mu$m (F115W) and 4.05$\mu$m (F405N), extinction ratios are
consistent across regions and align with values predicted 
by the \citep{Fritz2011} power law. Conversely, the extinction ratio at 3.23$\mu$m
(F323N) reveals substantial deviations, with NRCB3 exhibiting the highest ratio
of $4.18 \pm 0.43$, significantly larger than the other regions.
However the 3.23$mu$m calculated extinction ratios were in every region were
still significantly greater than the value predicted by the \citep{Fritz2011} power law. 
These significant
deviations likely arise from local physical conditions near Sagittarius A*,
including potential unevenly distributed ice absorption features. 

Our analysis underscores the need to account for spatial variations in
extinction when studying obscured environments like the Galactic Center. By
leveraging JWST's high-resolution data and advanced modeling techniques, we
provide new insights into the properties of interstellar dust and its impact on
extinction laws. Future work should aim to interpolate our discretized extinction law for all near to mid-infrared
wavelengths and explore the physical mechanisms
underlying these variations, such as the role of ice composition and dust grain
properties, using complementary spectroscopic and synthetic modeling techniques.
These efforts will refine our understanding of extinction laws and their broader
implications for the study of dense and obscured regions in the Milky Way.

\begin{figure*}
  \centering
    \includegraphics[height=\textheight]{nrcb3_plots.png}
    \caption{NRCB1 and NRCB3 examples of identifying well-defined RC stars from
        F115W - F212N vs. F115W CMDs and mapping them to higher wavelengths to
    determine if the extinction vectors align well. The red line indicates OCF's
calculated fit and the blue line indicates the predicted slope by the \citep{Fritz2011} power law.}
    \label{fig:nrcb3_plots}
\end{figure*}


\bibliographystyle{aasjournal}
\bibliography{draft}



\end{document}

